% Pregunta 148 (2022-julio, n.º 3): fuera del óptimo |RMS| > px/py; subida de px.
% A escala, con U = xy, R = 4, py = 1:
%   Izquierda: px = 1, recta x + y = 4. A = (1, 3), U = 3, pendiente de la curva en A = -3 (> 1 en valor absoluto).
%              E = (2, 2), U = 4: tangencia con pendiente -1.
%   Derecha: px pasa de 1 a 2, recta 2x + y = 4 (gira sobre (0, 4)). E' = (1, 2), U = 2 < 4.
% Autor: Víctor Gutiérrez Marcos
\documentclass[tikz,border=6pt]{standalone}
\input{tikz-tcee}
\begin{document}
\begin{tikzpicture}[tcee,x=1cm,y=1cm]
  % ---------- Panel izquierdo: opción a ----------
  \begin{scope}
    \node[etiqueta,anchor=south west,align=left] at (0,5.35) {a) misma renta y precios};
    \draw[eje] (0,0) -- (0,5.1) node[left] {$y$};
    \draw[eje] (0,0) -- (5.1,0) node[below] {$x$};
    \draw[curva1] (0,4) -- (4,0) node[etiqueta,below,text=tceemorado] {$R^0$};
    \draw[curva2,domain=0.64:4.7,samples=60] plot (\x,{3/\x}) node[etiqueta,right,text=tceeazul] {$U_A$};
    \draw[curva3,domain=0.86:4.8,samples=60] plot (\x,{4/\x}) node[etiqueta,above,text=tceeverde] {$U_E$};
    \draw[flecha] (1.2,2.95) -- (1.85,2.3);
    \node[punto] at (1,3) {}; \node[etiqueta,left] at (1,3) {$A$};
    \node[punto] at (2,2) {}; \node[etiqueta,above right] at (2,2) {$E$};
    \node[etiqueta,align=left,anchor=west] at (2.3,3.85) {$|RMS_A|>\frac{p_x}{p_y}$};
    \draw[guia] (2,0) node[below,text=black] {$x^*$} -- (2,2);
    \draw[guia] (1,0) node[below,text=black] {$x^0$} -- (1,3);
  \end{scope}
  % ---------- Panel derecho: opción c ----------
  \begin{scope}[xshift=6.6cm]
    \node[etiqueta,anchor=south west,align=left] at (0,5.35) {c) sube $p_x$};
    \fill[area2] (2,0) -- (4,0) -- (0,4) -- cycle;
    \draw[eje] (0,0) -- (0,5.1) node[left] {$y$};
    \draw[eje] (0,0) -- (5.1,0) node[below] {$x$};
    \draw[curva1] (0,4) -- (4,0) node[etiqueta,below,text=tceemorado] {$R^0$};
    \draw[curvanueva,desplazada] (0,4) -- (2,0) node[etiqueta,below,text=tceerojo] {$R^1$};
    \draw[curva3,domain=0.86:4.8,samples=60] plot (\x,{4/\x}) node[etiqueta,above,text=tceeverde] {$U_E$};
    \draw[curvanueva,domain=0.44:4.7,samples=60] plot (\x,{2/\x}) node[etiqueta,right,text=tceerojo] {$U_{E'}$};
    \node[punto] at (2,2) {}; \node[etiqueta,above right] at (2,2) {$E$};
    \node[punto,fill=tceerojo] at (1,2) {}; \node[etiqueta,left,text=tceerojo] at (1,2) {$E'$};
    \node[etiqueta,text=tceedorado!70!black,align=center] at (2.85,0.24) {se pierde};
  \end{scope}
\end{tikzpicture}
\end{document}
